% ============================================================================
% SEC03.TEX - Section 7.3: Memicu Ledakan via Script
% ============================================================================

\section{Memicu Ledakan via Script}

Animasi tidak berguna kalau tidak muncul saat kejadian. Kita harus memunculkan (Instantiate) efek ledakan ini saat asteroid hancur.

\subsection{Jadikan Prefab}
\begin{enumerate}
    \item Tarik \texttt{ExplosionVFX} dari Hierarchy ke folder \textbf{Prefabs}.
    \item Hapus \texttt{ExplosionVFX} dari Hierarchy.
\end{enumerate}

\subsection{Update AsteroidController.cs}
Tambahkan referensi prefab ledakan dan instansiasi sebelum destroy.

\begin{lstlisting}[language={[Sharp]C}]
public GameObject explosionPrefab; // Drag Prefab Explosion ke sini di Inspector

private void OnTriggerEnter2D(Collider2D other)
{
    if (other.CompareTag("Bullet"))
    {
        // Munculkan Ledakan
        Instantiate(explosionPrefab, transform.position, transform.rotation);
        
        Destroy(gameObject);
        Destroy(other.gameObject);
    }
}
\end{lstlisting}

\subsection{Auto Destroy Explosion}
Efek ledakan akan menumpuk di Hierarchy jika tidak dihapus. Tambahkan script sederhana \texttt{AutoDestroy.cs} pada prefab \texttt{ExplosionVFX}:

\begin{lstlisting}[language={[Sharp]C}]
public class AutoDestroy : MonoBehaviour
{
    void Start()
    {
        // Hancurkan diri sendiri setelah animasi selesai (misal 1 detik)
        Destroy(gameObject, 1f);
    }
}
\end{lstlisting}
